\documentclass[11pt]{article}

%%%%%%%%%%%%%%%%%%%%%%%%%%%%%%%% Algunos Paquetes Necesarios
\usepackage{fancyhdr, graphicx}
\usepackage[utf8]{inputenc}
\usepackage[T1]{fontenc}
\usepackage{lmodern}
\usepackage[spanish]{babel}
\usepackage{babelbib}
\usepackage[margin=0.85in]{geometry}
\usepackage{amssymb}
\usepackage{amsmath, amsthm, amsfonts}
\usepackage[table]{xcolor}
\usepackage{longtable}
\usepackage{enumitem}
\usepackage{multicol}
\decimalpoint
\usepackage{colortbl}
\usepackage{listings}
\usepackage{float}
\usepackage{placeins}
\usepackage{hyperref}

\definecolor{pastelblue}{RGB}{173,216,230}

\lstset{
    language=Python,
    basicstyle=\ttfamily\small,
    keywordstyle=\color{cyan!50!black},
    stringstyle=\color{orange!85!black},
    commentstyle=\color{teal!60!black},
    morecomment=[l][\color{magenta!60!black}]{\#},
    backgroundcolor=\color{white},
    showstringspaces=false,
    tabsize=4,
    frame=single,
    breaklines=true,
    captionpos=b
}

\hypersetup{
    colorlinks=true,
    linkcolor=blue!45!black,
    urlcolor=blue!45!black
}

\setlength{\emergencystretch}{2em}

%%%%%%%%%%%%%%%%%%%%%%%%%%%%%%%% Rutas de las imágenes
% Las rutas se interpretan desde la raiz del proyecto de Overleaf.
% Todo el material grafico vive dentro de imagenes/:
%   imagenes/Imagenes de prueba/  imagenes originales del proyecto
%   imagenes/Figuras/Fourier/     figuras de la Parte 1
%   imagenes/Figuras/FFT/         figuras de los incisos 1, 2 y 3
%   imagenes/Figuras/CNN/         figuras de los incisos 4 y 5
%   imagenes/Promts/              capturas de prompts y respuestas de IA
\graphicspath{{imagenes/}}

%%%%%%%%%%%%%%%%%%%%%%%%%%%%%%%% Capturas de prompts y respuestas
% \capturacnn{ruta}{descripcion}: inserta la imagen si existe;
% si todavia no existe, muestra un espacio reservado.
\newcommand{\capturacnn}[2]{%
  \par\medskip\noindent
  \IfFileExists{#1}{%
    \begin{center}
    \includegraphics[
        width=\linewidth,
        height=.28\textheight,
        keepaspectratio
    ]{#1}
    \end{center}%
  }{%
    \fbox{%
      \parbox[c][.15\textheight][c]{
        \dimexpr\linewidth-2\fboxsep-2\fboxrule\relax
      }{%
        \centering
        \textbf{#2}
        \par\smallskip
        Espacio para la captura real
        \par\smallskip
        {\small\path{#1}}
      }%
    }
  }%
  \par\smallskip
  \noindent{\small\textit{#2}}
  \par
}

% Evidencia de IA de Maria Claudia: \evidenciacnn{inciso}{numero}{texto}
% usa imagenes/Figuras/CNN/evidencia_N.jpg (prompt a la izquierda, respuesta a la derecha)
\newcommand{\evidenciacnn}[3]{%
  \subsubsection*{Registro de la consulta a IA: inciso #1}
  \IfFileExists{imagenes/Figuras/CNN/evidencia_#2.jpg}{%
    \begin{center}
    \includegraphics[width=\linewidth]{imagenes/Figuras/CNN/evidencia_#2.jpg}
    \end{center}\vspace{-2mm}%
    {\small\textit{Prompt (izquierda) y respuesta de la IA (derecha) del inciso #1.}}\par
  }{%
    \fbox{\parbox[c][.12\textheight][c]{.95\linewidth}{\centering Espacio para la captura del inciso #1}}\par
  }
  \paragraph{Uso de la asistencia.} #3
}

% Evidencia de IA de Antonio (un par de capturas por tema, lado a lado)
% usa imagenes/Promts/Tony/<clave>_prompt.png y <clave>_respuesta.png
\newcommand{\capturalado}[2]{%
  \parbox[t]{.485\linewidth}{\centering
  \IfFileExists{#1}{\includegraphics[width=\linewidth,height=.2\textheight,keepaspectratio]{#1}}{%
    \fbox{\parbox[c][.1\textheight][c]{.8\linewidth}{\centering\small Espacio para la captura real\par\path{#1}}}}\par
  {\small\textit{#2}}}%
}
\newcommand{\evidenciaantonio}[3]{%
  \paragraph{#1.}
  \noindent\capturalado{imagenes/Promts/Tony/#2_prompt.png}{Prompt}\hfill
  \capturalado{imagenes/Promts/Tony/#2_respuesta.png}{Respuesta de la IA}\par\smallskip
  #3\par\medskip
}

% Hoja de capturas de la conversacion de Christian (apendice)
\newcommand{\hojachristian}[1]{%
  \begin{center}\includegraphics[width=\linewidth,height=.88\textheight,keepaspectratio]{imagenes/Figuras/Apendice/hoja_#1.jpg}\end{center}%
}

% Interfaz del notebook: dos capturas (superior e inferior) apiladas, una pagina por caso
% \interfazpar{carpeta}{N}{etiqueta}{leyenda}  usa imagenes/interfaz/<carpeta>/superior_N_<N>.png e inferior_N_<N>.png
\newcommand{\interfazpar}[4]{%
  \begin{figure}[p]
  \centering
  \includegraphics[width=.58\linewidth,height=.54\textheight,keepaspectratio]{imagenes/interfaz/#1/superior_N_#2.png}\par\smallskip
  \includegraphics[width=.58\linewidth,height=.4\textheight,keepaspectratio]{imagenes/interfaz/#1/inferior_N_#2.png}
  \caption{#4}
  \label{#3}
  \end{figure}
}

% Figura simple a ancho completo
\newcommand{\figura}[3]{%
  \begin{figure}[!htbp]
  \centering
  \includegraphics[width=.76\linewidth]{#1}
  \caption{#3}
  \label{#2}
  \end{figure}
}

% Imagen que puede faltar (por ejemplo una captura de pantalla pendiente)
\newcommand{\figuraopcional}[4]{%
  \begin{figure}[H]
  \centering
  \IfFileExists{#1}{%
    \includegraphics[width=#4\linewidth]{#1}%
  }{%
    \fbox{\parbox[c][.18\textheight][c]{.9\linewidth}{\centering
      \textbf{Espacio para la captura real}\par\smallskip
      {\small\path{#1}}}}%
  }
  \caption{#3}
  \label{#2}
  \end{figure}
}

%%%%%%%%%%%%%%%%%%%%%%%%%%%%%%%% Logo de la universidad
\newcommand{\logouniversidad}{%
  \begin{center}
  \IfFileExists{UGlogo_HR.png}{%
    \includegraphics[width=.3\textwidth]{UGlogo_HR.png}
  }{%
    \IfFileExists{UGlogo_HR.pdf}{%
      \includegraphics[width=.3\textwidth]{UGlogo_HR.pdf}
    }{%
      \IfFileExists{UGlogo_HR.jpg}{%
        \includegraphics[width=.3\textwidth]{UGlogo_HR.jpg}
      }{%
        \fbox{%
          \parbox[c][1.3cm][c]{.3\textwidth}{%
            \centering
            \small
            Logo de la universidad\\
            \texttt{UGlogo\_HR}
          }%
        }
      }
    }
  }
  \end{center}
}

%%%%%%%%%%%%%%%%%%%%%%%%%%%%%%%% Figuras de resultados
\newcommand{\panelcnn}[3]{%
  \includegraphics[width=\linewidth]{imagenes/Figuras/CNN/#1/#1_img#2_#3.jpg}%
}

% Cuatro comparaciones de dos paneles.
\newcommand{\galeriacnn}[3]{%
  \begin{figure}[!htbp]
  \centering

  \begin{minipage}{.485\linewidth}
    \centering
    \panelcnn{#1}{1}{#2}
  \end{minipage}
  \hfill
  \begin{minipage}{.485\linewidth}
    \centering
    \panelcnn{#1}{2}{#2}
  \end{minipage}

  \par\medskip

  \begin{minipage}{.485\linewidth}
    \centering
    \panelcnn{#1}{3}{#2}
  \end{minipage}
  \hfill
  \begin{minipage}{.485\linewidth}
    \centering
    \panelcnn{#1}{4}{#2}
  \end{minipage}

  \caption{#3}
  \label{fig:cnn-#1}
  \end{figure}
}

% Comparaciones con varios paneles a ancho completo.
\newcommand{\comparacionescnn}[3]{%
  \begin{figure}[!htbp]
  \centering

  \panelcnn{#1}{1}{#2}
  \par\smallskip
  \panelcnn{#1}{2}{#2}
  \par\smallskip
  \panelcnn{#1}{3}{#2}
  \par\smallskip
  \panelcnn{#1}{4}{#2}

  \caption{#3}
  \label{fig:cnn-#1}
  \end{figure}
}

% Restauracion con DnCNN: cuatro imagenes en una figura.
\newcommand{\comparacionesdncnn}{%
  \begin{figure}[!htbp]
  \centering
  \panelcnn{4d}{1}{dncnn_comparacion}\par\smallskip
  \panelcnn{4d}{2}{dncnn_comparacion}\par\smallskip
  \panelcnn{4d}{3}{dncnn_comparacion}\par\smallskip
  \panelcnn{4d}{4}{dncnn_comparacion}
  \caption{Restauración con DnCNN de img1 a img4 (una por fila). Cada fila muestra la referencia, la entrada con ruido de 4b y la salida restaurada, con el mismo orden de paneles y la misma escala de intensidades.}
  \label{fig:cnn-4d}
  \end{figure}
}
\begin{document}

%%%%%%%%%%%%%%%%%%%%%%%%%%%%%%%%% Titulo
{\color{white}.}
\vspace{3 in}

\begin{center}
\Huge
\textsc{MAP $-$ FFT y DFT}
\end{center}

\vspace{3 in}
\hrule height 1.5pt
\vspace{0.1 in}

%%%%%%%%%%%%%%%%%%%%%%%%%%%%%%%%% Datos Grupo
{\large
\noindent\textbf{Integrantes:}

\begin{itemize}[nolistsep]
    \item María Claudia Lainfiesta Herrera $-$ $24000149$.
    \item Antonio Papadopolo Castellanos $-$ $24002045$.
    \item Christian Alejandro Vasquez Paz $-$ $24001777$.
\end{itemize}
}

\vspace{0.2 in}
\hrule height 1.5pt

\newpage
\vspace{0.1 in}



\begin{center}
\Large
\textbf{Guía Práctica}
\end{center}

%%%%%%%%%%%%%%%%%%%%%%%%%%%%%%%%%%%%%%%%%%%%%%%%%%%%%%
% INICIO DEL APARTADO DE CHRISTIAN ALEJANDRO VASQUEZ PAZ
\subsection*{Parte 1 - Series de Fourier}

\noindent

\paragraph{Convenciones.}
Las dos funciones del enunciado tienen período $T=4$ y se integran en
$[-2,2]$, de modo que $\omega_0=2\pi/T=\pi/2$. La serie truncada al armónico $N$ es
\begin{equation}
 S_N(t)=a_0+\sum_{n=1}^{N}\bigl[a_n\cos(n\omega_0 t)+b_n\sin(n\omega_0 t)\bigr],
 \label{eq:sn}
\end{equation}
con $a_0=\frac1T\int f\,dt$, $a_n=\frac2T\int f\cos(n\omega_0t)\,dt$ y
$b_n=\frac2T\int f\sin(n\omega_0t)\,dt$. El término constante es $a_0$, sin dividirlo otra vez entre dos.
La expresión (\ref{eq:sn}) tiene $2N+1$ términos: $a_0$, los $N$ cosenos y los $N$ senos
(algunos coeficientes valen cero). La energía y la integral cuadrada del error son
\begin{equation}
 E_f=\int_{-2}^{2}f^2(t)\,dt,\qquad
 \mathrm{ICE}(N)=E_f-\Bigl\{a_0^2T+\frac T2\sum_{n=1}^{N}(a_n^2+b_n^2)\Bigr\},
\end{equation}
y el criterio de aceptación es $\mathrm{ICE}(N)\le0.02\,E_f$.

\paragraph{Coeficientes en forma cerrada.}
Al integrar a mano se obtienen, con $k=n\pi/2$:
\begin{align*}
 f_1:&\quad a_0=\tfrac12,\quad a_n=\frac{2\sin(n\pi/2)}{n\pi},\quad b_n=0,\\
 f_2:&\quad a_0=1,\quad a_n=\frac{4\sin(n\pi/2)}{n\pi},\quad
 b_n=\frac{\sin k}{k^2}-\frac{\cos k}{k},
\end{align*}
con energías $E_{f_1}=2$ y $E_{f_2}=26/3$. El código calcula los coeficientes por integración numérica
(\texttt{scipy.integrate.quad}, separando los puntos de discontinuidad $-1$, $0$ y $1$) y las fórmulas anteriores se usaron
para verificarlos hasta $n=11$ con tolerancia $10^{-9}$.

%%%%%%%%%%%%%%%%%%%%%%%%%%%%%%%%%%%%%%%%%%%%%%%%%%%%%%
\subsubsection*{Inciso 1. Interfaz gráfica}

\paragraph{Instrucción.}
Implementar una interfaz gráfica que permita seleccionar la función ($f_1$ o $f_2$) y que muestre
el período, los coeficientes de la serie, la gráfica aproximada con $2N+1$ términos, el valor de $N$
y la ecuación de la serie trigonométrica.

\paragraph{Respuesta y explicación del código.}
La interfaz se construyó con \texttt{ipywidgets} dentro del notebook, y se actualiza sola al mover cualquier
control (\texttt{widgets.interactive\_output}). Cada dato pedido en el enunciado tiene su propio elemento:
\begin{itemize}[nolistsep]
 \item \textbf{Selección de función:} menú desplegable con $f_1$ (pulso rectangular), $f_2$ (rampa por tramos)
 y la señal propuesta con IA del inciso 4.
 \item \textbf{Período:} tarjeta con $T=4$ y la frecuencia angular $\omega_0=\pi/2$.
 \item \textbf{Valor de $N$:} deslizador de $0$ a $30$ armónicos. El botón <<Calcular $N$ mínimo (2\,\%)>>
 busca el menor $N$ con $\mathrm{ICE}(N)\le0.02E_f$ y el botón <<Caso IA: $N=15$>> prepara el inciso 5.
 \item \textbf{Coeficientes:} $a_0$ y una tabla con $a_n$ y $b_n$ para $n=1,\dots,N$.
 \item \textbf{Gráfica:} señal original (extendida periódicamente, de 1 a 5 períodos) frente a la serie
 truncada; incluye opción para mostrar u ocultar la cuadrícula.
 \item \textbf{Ecuación:} expresión $S_N(t)$ con los coeficientes redondeados a seis decimales
 (se omiten términos menores que $10^{-10}$ solo al escribirla; la gráfica usa los valores sin redondear).
 \item \textbf{Calidad de la aproximación:} tarjeta con energía, ICE, porcentaje $\mathrm{ICE}/E_f$ y si cumple el 2\,\%.
\end{itemize}
Las funciones \texttt{calcular\_a0}, \texttt{calcular\_an} y \texttt{calcular\_bn} integran los coeficientes;
\texttt{serie\_fourier} suma la serie truncada; \texttt{calcular\_energia} y \texttt{calcular\_ice} dan $E_f$ y la ICE,
y \texttt{encontrar\_n} busca el menor $N$ que cumple el 2\,\%:

\begin{lstlisting}
def calcular_an(funcion, n):
    def integrando(t):
        return funcion(t) * np.cos(n * omega_0 * t)
    integral, error_estimado = quad(integrando, -2, 2, points=[-1, 0, 1])
    return (2 / T) * integral

def calcular_ice(funcion, N):
    energia = calcular_energia(funcion)
    a0 = calcular_a0(funcion)
    suma = sum(calcular_an(funcion, n)**2 + calcular_bn(funcion, n)**2
               for n in range(1, N + 1))
    return energia - (T * a0**2 + (T / 2) * suma)
\end{lstlisting}

Las capturas de la interfaz funcionando, con todos los datos que pide el enunciado, se muestran en el inciso 2
(figuras \ref{fig:int-f1-3} a \ref{fig:int-fia-15}).

%%%%%%%%%%%%%%%%%%%%%%%%%%%%%%%%%%%%%%%%%%%%%%%%%%%%%%
\subsubsection*{Inciso 2. Documentación para $N=3,5,10$}

\paragraph{Instrucción.}
Documentar paso a paso, con todos los datos generados por la interfaz, los valores $N=3,5,10$ para cada función.

\paragraph{Respuesta.}
Cada una de las siguientes figuras es la interfaz del notebook para un valor de $N$, en dos capturas (parte superior arriba y
parte inferior abajo). En ellas se ven, para cada caso: la función seleccionada, el período $T=4$ y $\omega_0=\pi/2$, el valor
de $N$, la energía y la ICE con su porcentaje, la gráfica de la señal original frente a la serie truncada, la ecuación
$S_N(t)$ con sus $2N+1$ posiciones y la tabla de coeficientes $a_0$, $a_n$ y $b_n$. La tabla \ref{tab:resumen6} resume la
energía y la ICE de los seis casos, y las figuras \ref{fig:int-f1-3} a \ref{fig:int-f1-10} corresponden a $f_1$ y
\ref{fig:int-f2-3} a \ref{fig:int-f2-10} a $f_2$.

\begin{table}[!htbp]
\centering
\caption{Energía, ICE y cumplimiento del criterio $\mathrm{ICE}\le0.02E_f$ en los seis casos.}
\label{tab:resumen6}
\begin{tabular}{llrrrrl}
\hline
\rowcolor{pastelblue}
Función & $N$ & $2N+1$ & $E_f$ & ICE & ICE$/E_f$ (\%) & ¿Cumple 2\,\%?\\
\hline
f1 & 3 & 7 & 2.000000 & 0.099367 & 4.9684 & No\\
f1 & 5 & 11 & 2.000000 & 0.066944 & 3.3472 & No\\
f1 & 10 & 21 & 2.000000 & 0.040395 & 2.0198 & No\\
\hline
f2 & 3 & 7 & 8.666667 & 0.528926 & 6.1030 & No\\
f2 & 5 & 11 & 8.666667 & 0.348049 & 4.0159 & No\\
f2 & 10 & 21 & 8.666667 & 0.198378 & 2.2890 & No\\
\hline
\end{tabular}
\end{table}

\interfazpar{f1}{3}{fig:int-f1-3}{Interfaz para $f_1$ con $N=3$.}
\interfazpar{f1}{5}{fig:int-f1-5}{Interfaz para $f_1$ con $N=5$.}
\interfazpar{f1}{10}{fig:int-f1-10}{Interfaz para $f_1$ con $N=10$.}
\interfazpar{f2}{3}{fig:int-f2-3}{Interfaz para $f_2$ con $N=3$.}
\interfazpar{f2}{5}{fig:int-f2-5}{Interfaz para $f_2$ con $N=5$.}
\interfazpar{f2}{10}{fig:int-f2-10}{Interfaz para $f_2$ con $N=10$.}
\FloatBarrier


%%%%%%%%%%%%%%%%%%%%%%%%%%%%%%%%%%%%%%%%%%%%%%%%%%%%%%
\subsubsection*{Inciso 3. Preguntas de análisis}

\paragraph{a) ¿En cuál de las dos funciones converge más rápido la serie?}
En los casos estudiados, \textbf{$f_1$ presenta menor error relativo}: para $N=3,5,10$ obtiene
4.9684\,\%, 3.3472\,\% y 2.0198\,\%, frente a 6.1030\,\%, 4.0159\,\% y 2.2890\,\% de $f_2$.
Con el mismo número de armónicos, $f_1$ ofrece entonces una mejor aproximación según el criterio energético relativo.
Esto es coherente con los saltos: $f_1$ tiene dos saltos de magnitud $1$, mientras que $f_2$ tiene saltos de
magnitudes $1$ y $3$.

No obstante, \textbf{ambas alcanzan el umbral del 2\,\% por primera vez en $N=11$}, así que no debe concluirse que
$f_1$ necesite menos armónicos para ese umbral particular. Los coeficientes de ambas decaen como $1/n$, y la energía del
error decae como $1/N$; la diferencia observada es de constantes, no de orden de convergencia.

\paragraph{b) ¿Es útil calcular la ICE para ambas funciones?}
Sí. Permite cuantificar el error cuadrático sobre un período y elegir $N$ con un criterio reproducible:
\[
 \mathrm{ICE}(N)=\int_{-2}^{2}[f(t)-S_N(t)]^2dt=E_f-Ta_0^2-\frac T2\sum_{n=1}^N(a_n^2+b_n^2),
\]
igualdad que se debe a la ortogonalidad de las funciones trigonométricas. Como $E_{f_1}=2$ y $E_{f_2}=26/3$ son distintas,
conviene comparar $\mathrm{ICE}/E_f$ y no el error absoluto.

La ICE, sin embargo, no mide el error máximo en cada punto. Cerca de los saltos persisten las oscilaciones de Gibbs:
al aumentar $N$ la zona afectada se estrecha, pero el sobreimpulso relativo no desaparece, y en el salto la serie converge
al promedio de los límites laterales. Cumplir el 2\,\% de energía no significa estar a menos del 2\,\% de la función en
todos los instantes.

\paragraph{c) ¿Con cuántas iteraciones se obtiene una mejor aproximación?}
Interpretamos <<iteraciones>> como el armónico máximo $N$ de la suma. Entre los tres casos pedidos, $N=10$ da la menor ICE
en cada función, pero todavía no cumple el 2\,\%. El mínimo que lo cumple es $N=11$ para ambas (figura \ref{fig:ice}):

\begin{table}[!htbp]
\centering
\caption{Error relativo $\mathrm{ICE}/E_f$ (\%) alrededor del primer $N$ que cumple el criterio.}
\label{tab:n11}
\begin{tabular}{lrr}
\hline
\rowcolor{pastelblue}
Función & $N=10$ & $N=11$\\
\hline
$f_1$ & 2.0198 & 1.6848\\
$f_2$ & 2.2890 & 1.9795\\
\hline
\end{tabular}
\end{table}

Agregar un armónico reduce la ICE en $\frac T2(a_n^2+b_n^2)\ge0$, por lo que la ICE nunca aumenta; puede quedar igual cuando los
nuevos coeficientes son cero, como ocurre con los armónicos pares de $f_1$. Ningún $N$ finito reproduce exactamente las
discontinuidades.

\figura{imagenes/Figuras/Fourier/ice_vs_N.jpg}{fig:ice}{
 Error relativo $\mathrm{ICE}(N)/E_f$ para $f_1$ y $f_2$. La línea roja es el criterio del 2\,\% y la vertical punteada marca $N=11$.}

%%%%%%%%%%%%%%%%%%%%%%%%%%%%%%%%%%%%%%%%%%%%%%%%%%%%%%
\subsubsection*{Incisos 4 a 6. Señal periódica generada con IA}

\paragraph{Instrucciones.}
(4) Pedir a un LLM el código de una función periódica original y no trivial; (5) analizarla con la interfaz,
documentar su serie para $N=15$ y juzgar si la aproximación es adecuada; (6) graficar la señal y su aproximación.

\paragraph{Señal propuesta.}
Con asistencia de IA se propuso una señal sintética que combina tres componentes sinusoidales y una conmutación brusca
(es un ejemplo matemático, no un modelo validado de un instrumento real). En el intervalo base $[-2,2)$:
\begin{equation}
 g(t)=\sin\frac{\pi t}{2}+0.45\sin\frac{3\pi t}{2}+0.20\cos\frac{5\pi t}{2}+0.35\,q(t),\qquad
 q(t)=\begin{cases}-1,&-2\le t<0,\\ +1,&0\le t<2,\end{cases}
\end{equation}
extendida con período $T=4$. Las sinusoides son los armónicos $1$, $3$ y $5$, y $q$ introduce saltos en $t=0$ y en los
extremos del período, de modo que se necesitan infinitos armónicos para representarla exactamente. Las consultas
al asistente se documentan al final de esta parte (página \pageref{anexo:ia}).

\begin{lstlisting}
def funcion_ia(t):
    t = np.asarray(t, dtype=float)
    cuadrada = np.where(t < 0, -1.0, 1.0)
    return (np.sin(omega_0 * t)
            + 0.45 * np.sin(3 * omega_0 * t)
            + 0.20 * np.cos(5 * omega_0 * t)
            + 0.35 * cuadrada)
\end{lstlisting}

\paragraph{Verificación independiente de los coeficientes.}
Por linealidad y por la serie de la onda cuadrada, $a_0=0$, $a_n=0.20$ solo si $n=5$ (y cero en los demás), y
\[
 b_n=\mathbf 1_{n=1}+0.45\,\mathbf 1_{n=3}+\begin{cases}\dfrac{1.4}{n\pi},&n\text{ impar},\\[2mm]0,&n\text{ par},\end{cases}
\]
donde $\mathbf 1_{\cdot}$ vale $1$ si se cumple la condición y $0$ en otro caso. Estas expresiones coinciden con los
coeficientes calculados por integración (con tolerancia $10^{-9}$) para $n=1,\dots,15$. Es una comprobación independiente
de la implementación.

\paragraph{Serie truncada para $N=15$ en la interfaz.}
La figura \ref{fig:int-fia-15} muestra la interfaz con la señal propuesta por la IA y $N=15$: período $T=4$, $\omega_0=\pi/2$,
$2N+1=31$ posiciones, energía e ICE, gráfica de la señal original y de la aproximación, ecuación $S_{15}(t)$ y tabla de
coeficientes ($a_0=0$, $a_5=0.20$ y los $b_n$ de los armónicos impares).

\interfazpar{fIA}{15}{fig:int-fia-15}{Interfaz para la señal propuesta con IA con $N=15$.}
\FloatBarrier

\paragraph{¿Es adecuada la aproximación con $N=15$?}
Sí, \textbf{según el criterio de energía}. La energía es $E_g\approx5.024916$ y $\mathrm{ICE}(15)\approx0.012396$, es decir
$\mathrm{ICE}/E_g\approx0.2467$\,\%, muy por debajo del 2\,\%. De hecho el criterio ya se cumple desde $N=5$
(0.6528\,\%). La energía total tiene el valor analítico
\[
 E_g=2(1+0.45^2+0.20^2)+4(0.35)^2+\frac{16(0.35)}{\pi}\Bigl(1+\frac{0.45}{3}\Bigr)\approx 5.024916,
\]
que coincide con la integración numérica y no es la suma simple de las energías de los términos, porque la onda cuadrada
comparte armónicos con los senos de frecuencias $1$ y $3$ y aparecen términos cruzados.

La conclusión es energética y global, no afirma precisión uniforme en cada instante. Como una suma finita de senos y
cosenos es continua, cerca de $t=0$ y de los extremos del período quedan oscilaciones y una transición suave en lugar del
salto (figura \ref{fig:ia}, panel inferior).

\figura{imagenes/Figuras/Fourier/ia_N15.jpg}{fig:ia}{
 Señal propuesta con IA y su serie truncada con $N=15$ sobre dos períodos (arriba) y detalle cerca del salto en $t=0$ (abajo).}

\clearpage
\subsubsection*{Registro de las consultas a IA (Parte 1)}
\label{anexo:ia}

Las siguientes hojas reúnen las 51 capturas de la conversación continua con el asistente Codex que se utilizó para desarrollar la
Parte 1: se avanzó paso a paso con las funciones $f_1$ y $f_2$, la integración numérica de los coeficientes, la energía, la ICE,
la interfaz y la señal periódica propuesta. Están en orden (capturas 1 a 51) y reducidas para el reporte; las capturas originales
están en \path{imagenes/Promts/Christian/}. Los valores numéricos del reporte provienen de la ejecución del notebook, no de cifras propuestas por la IA.

\hojachristian{1}
\clearpage
\hojachristian{2}
\clearpage
\hojachristian{3}

% FIN DEL APARTADO DE CHRISTIAN ALEJANDRO VASQUEZ PAZ

\clearpage
%%%%%%%%%%%%%%%%%%%%%%%%%%%%%%%%%%%%%%%%%%%%%%%%%%%%%%
% INICIO DEL APARTADO DE ANTONIO PAPADOPOLO CASTELLANOS
\subsection*{Parte 2 - FFT y filtros de imagen avanzados}

\noindent

%%%%%%%%%%%%%%%%%%%%%%%%%%%%%%%%%%%%%%%%%%%%%%%%%%%%%%
\subsubsection*{Implementación del algoritmo FFT de Cooley-Tukey}

\paragraph{Instrucción.}
Implementar el algoritmo FFT recursivo del enunciado y usarlo en los incisos 1, 2 y 3.

\paragraph{Respuesta y explicación del código.}
El algoritmo se implementó sin usar \texttt{np.fft.fft} ni \texttt{np.fft.fft2} para calcular transformadas.
Sigue el pseudocódigo del enunciado: caso base $N=1$; transformadas recursivas de los índices pares e impares;
y combinación con $T=x_{\text{impar}}[k]\,e^{-2\pi ik/N}$,
\begin{equation}
 X[k]=X_{\text{par}}[k]+T,\qquad X[k+N/2]=X_{\text{par}}[k]-T,\qquad k=0,\dots,\tfrac N2-1.
\end{equation}
Al dividir el problema a la mitad en cada nivel, el costo pasa de $O(N^2)$ de la DFT directa a $O(N\log N)$,
y por la misma razón la longitud debe ser potencia de dos; el código lanza un error claro en otro caso.

\begin{lstlisting}
def fft_recursiva(x):
    x = np.asarray(x, dtype=complex)
    N = x.shape[-1]
    if not es_potencia_de_2(N):
        raise ValueError(f"La longitud debe ser potencia de 2, se recibio N={N}")
    if N == 1:                              # caso base
        return x
    x_par = fft_recursiva(x[..., 0::2])     # indices pares
    x_impar = fft_recursiva(x[..., 1::2])   # indices impares
    k = np.arange(N // 2)
    T = x_impar * np.exp(-2j * np.pi * k / N)
    return np.concatenate([x_par + T, x_par - T], axis=-1)
\end{lstlisting}

La inversa usa la identidad $\mathrm{IFFT}(X)=\overline{\mathrm{FFT}(\overline{X})}/N$. Para imágenes se aplica la FFT
a todas las filas y luego a todas las columnas (FFT 2D separable), y la inversa 2D es análoga. El único uso de
\texttt{np.fft} en el proyecto es \texttt{fftshift}, que solo reordena los cuadrantes para llevar la frecuencia cero al
centro del espectro, sin calcular ninguna transformada.

\paragraph{Validación.}
La FFT propia se comparó con \texttt{np.fft} para longitudes $1,2,8,64,1024$ y para imágenes de
$64\times64$, $32\times128$ y $256\times256$; la inversa devuelve la señal original, y la DFT directa
(la fórmula $\mathbf X[k]=\sum_n x[n]e^{-2\pi ikn/N}$ del enunciado) coincide con la FFT. Todas las pruebas pasaron.

\paragraph{Preparación de las imágenes.}
Cada imagen se convierte a escala de grises, se redimensiona a $256\times256$ píxeles con interpolación Lanczos y se
normaliza a $[0,1]$. Las dimensiones deben ser potencias de dos porque la FFT recursiva lo exige; $256=2^8$ conserva
el detalle suficiente con un costo moderado.

\figura{imagenes/Figuras/FFT/originales_256.jpg}{fig:originales}{
 Las cuatro imágenes del proyecto ya preparadas ($256\times256$, escala de grises).}

\paragraph{Convención de frecuencia de corte.}
Se usa el radio normalizado $r$ medido desde el centro del espectro centrado: $r=0$ es la frecuencia cero y
$r=1$ es la esquina, la frecuencia más alta. Para una imagen de $H\times W$,
\begin{equation}
 r(u,v)=\frac{1}{\sqrt2}\sqrt{\Bigl(\frac{u-H/2}{H/2}\Bigr)^2+\Bigl(\frac{v-W/2}{W/2}\Bigr)^2}.
\end{equation}
Las máscaras ideales son
\begin{equation}
 H_{\text{bajas}}(u,v)=\begin{cases}1,&r\le f_c,\\0,&r>f_c,\end{cases}\qquad
 H_{\text{altas}}(u,v)=1-H_{\text{bajas}}(u,v).
\end{equation}
El filtrado calcula $F=\operatorname{centrar}(\mathrm{FFT2}(\text{img}))$, multiplica por la máscara y reconstruye
$g=\operatorname{Re}\bigl(\mathrm{IFFT2}(\operatorname{descentrar}(F\cdot H))\bigr)$. La parte imaginaria es solo error
numérico. En el pasa bajas el resultado se recorta a $[0,1]$; en el pasa altas no se recorta, porque el resultado tiene
media cero y valores negativos.

\figura{imagenes/Figuras/FFT/mascaras.jpg}{fig:mascaras}{
 Máscara pasa bajas ($f_c=0.22$), máscara pasa altas ($f_c=0.23$) y el radio normalizado $r$ (de 0 en el centro a 1 en la esquina).}



\FloatBarrier
%%%%%%%%%%%%%%%%%%%%%%%%%%%%%%%%%%%%%%%%%%%%%%%%%%%%%%
\subsubsection*{Inciso 1. Filtro pasa bajas ($f_c=0.22,\ 0.15,\ 0.96$)}

\paragraph{Instrucción.}
Aplicar el algoritmo a las imágenes con un filtro pasa bajas con $f_c=0.22$, $0.15$ y $0.96$, y explicar la diferencia entre las
imágenes originales y las generadas, justificando la respuesta con el análisis del filtro.

\paragraph{Respuesta y explicación del código.}
Se aplicó el filtro pasa bajas ideal a las cuatro imágenes con la FFT propia.

\begin{lstlisting}
#pasa bajas ideal con la FFT propia
for fc in [0.22, 0.15, 0.96]:
    filtrada = filtro_pasabajas(img, fc, tipo="ideal")
\end{lstlisting}

La figura \ref{fig:bajas} muestra los resultados para las cuatro imágenes y la figura \ref{fig:bajasdet} detalla img1 con la diferencia y el espectro filtrado (escala logarítmica); el comportamiento es el mismo en las demás. La tabla \ref{tab:bajas} da los números.
La energía conservada es $\sum|F\cdot H|^2/\sum|F|^2$, y el PSNR se calcula contra la imagen original
($\mathrm{PSNR}=10\log_{10}(1/\mathrm{MSE})$).

\begin{table}[!htbp]
\centering
\caption{Pasa bajas ideal: porcentaje de puntos del espectro que se conservan, energía conservada y PSNR contra la original.}
\label{tab:bajas}
\begin{tabular}{llrrr}
\hline
\rowcolor{pastelblue}
Imagen & $f_c$ & Puntos (\%) & Energía (\%) & PSNR (dB)\\
\hline
img1 & 0.22 & 7.59 & 98.41 & 23.44\\
img1 & 0.15 & 3.53 & 97.40 & 21.28\\
img1 & 0.96 & 99.66 & 100.00 & 71.98\\
\hline
img2 & 0.22 & 7.59 & 99.16 & 25.68\\
img2 & 0.15 & 3.53 & 98.70 & 23.76\\
img2 & 0.96 & 99.66 & 100.00 & 64.41\\
\hline
img3 & 0.22 & 7.59 & 99.35 & 26.55\\
img3 & 0.15 & 3.53 & 98.89 & 24.22\\
img3 & 0.96 & 99.66 & 100.00 & 69.41\\
\hline
img4 & 0.22 & 7.59 & 98.53 & 25.21\\
img4 & 0.15 & 3.53 & 97.72 & 23.31\\
img4 & 0.96 & 99.66 & 100.00 & 68.82\\
\hline
\end{tabular}
\end{table}

\figura{imagenes/Figuras/FFT/bajas_resumen.jpg}{fig:bajas}{Pasa bajas aplicado a las cuatro imágenes (una por fila): original y resultados con $f_c=0.22$, $0.15$ y $0.96$.}
\figura{imagenes/Figuras/FFT/img1_pasa_bajas.jpg}{fig:bajasdet}{Detalle para img1: por fila $f_c=0.22$, $0.15$ y $0.96$; columnas: original, filtrada, $|\text{original}-\text{filtrada}|$ y espectro filtrado (escala logarítmica).}

\paragraph{¿Cuál es la diferencia entre las imágenes originales y las generadas?}
Las imágenes generadas con $f_c=0.22$ y $f_c=0.15$ son versiones suavizadas de las originales: se conservan las formas
generales y las zonas de brillo uniforme, pero desaparecen los contornos nítidos y los detalles finos, y junto a los
bordes aparecen ondulaciones. Mientras menor es la frecuencia de corte, mayor es la pérdida: con $0.15$ el resultado se
ve más borroso que con $0.22$ (PSNR de 21.3 a 24.2 dB frente a 23.4 a 26.5 dB). En cambio, con $f_c=0.96$ la imagen generada
es prácticamente idéntica a la original (PSNR de 64 a 72 dB según la imagen), porque casi no se descarta contenido.

\paragraph{Justificación con el filtro pasa bajas.}
La máscara vale 1 dentro del radio $f_c$ y 0 fuera. Las frecuencias bajas describen variaciones lentas de intensidad,
como fondos y superficies, mientras que los bordes y las texturas son transiciones abruptas que necesitan muchas frecuencias
altas para representarse. Al eliminarlas, cada borde se reparte en una transición gradual, que se percibe como desenfoque;
el mapa de diferencia confirma que el error se concentra en los contornos.

Las cifras muestran por qué la imagen sigue siendo reconocible: con $f_c=0.22$ se conserva solo el 7.6\,\% de los puntos del
espectro y con $f_c=0.15$ apenas el 3.5\,\%, pero esos puntos retienen entre 98.4\,\% y 99.4\,\% de la energía (0.22) y entre
97.4\,\% y 98.9\,\% (0.15), porque en una imagen natural la energía se concentra en las bajas frecuencias. Con $f_c=0.96$ solo
se descarta el 0.3\,\% del espectro (las esquinas) y la energía conservada es prácticamente 100\,\%.

Las ondulaciones cercanas a los bordes se deben a que el corte de la máscara es abrupto. Multiplicar por una máscara de borde
nítido en frecuencia equivale a convolucionar con una función tipo $\operatorname{sinc}$ en el dominio espacial, cuyos
lóbulos laterales generan oscilaciones junto a cada discontinuidad. Es el fenómeno de Gibbs que ya aparece al truncar una serie de
Fourier: aproximar una discontinuidad con un número finito de armónicos produce sobreimpulsos cerca del salto.



\FloatBarrier
%%%%%%%%%%%%%%%%%%%%%%%%%%%%%%%%%%%%%%%%%%%%%%%%%%%%%%
\subsubsection*{Inciso 2. Filtro pasa altas ($f_c=0.23,\ 0.11,\ 0.91$)}

\paragraph{Instrucción.}
Aplicar el algoritmo a las imágenes con un filtro pasa altas con $f_c=0.23$, $0.11$ y $0.91$, y explicar la diferencia entre las
imágenes originales y las generadas, justificando la respuesta con el análisis del filtro.

\paragraph{Respuesta y explicación del código.}
El pasa altas usa la máscara complementaria $1-H_{\text{bajas}}$, que elimina la componente de frecuencia cero y deja
pasar lo que está fuera del radio $f_c$.

\begin{lstlisting}
#pasa altas ideal con la FFT propia
for fc in [0.23, 0.11, 0.91]:
    filtrada = filtro_pasaaltas(img, fc, tipo="ideal")
\end{lstlisting}

Como el resultado tiene media cero, para mostrarlo se escala a $[0,1]$ (el gris medio corresponde al valor cero); esa
normalización es solo visual. No se calcula PSNR, porque el resultado no debe parecerse a la imagen original. Las
figuras \ref{fig:altas} y \ref{fig:altasdet} y la tabla \ref{tab:altas} reúnen los resultados.

\begin{table}[!htbp]
\centering
\caption{Pasa altas ideal: puntos del espectro conservados, energía conservada y desviación estándar de la imagen resultante (escala 0 a 1).}
\label{tab:altas}
\begin{tabular}{llrrr}
\hline
\rowcolor{pastelblue}
Imagen & $f_c$ & Puntos (\%) & Energía (\%) & Desv. estándar\\
\hline
img1 & 0.23 & 91.69 & 1.4762 & 0.0656\\
img1 & 0.11 & 98.11 & 3.7700 & 0.1048\\
img1 & 0.91 & 1.70 & 0.0002 & 0.0007\\
\hline
img2 & 0.23 & 91.69 & 0.7930 & 0.0506\\
img2 & 0.11 & 98.11 & 1.7732 & 0.0757\\
img2 & 0.91 & 1.70 & 0.0005 & 0.0013\\
\hline
img3 & 0.23 & 91.69 & 0.6021 & 0.0454\\
img3 & 0.11 & 98.11 & 1.6797 & 0.0759\\
img3 & 0.91 & 1.70 & 0.0002 & 0.0008\\
\hline
img4 & 0.23 & 91.69 & 1.3931 & 0.0534\\
img4 & 0.11 & 98.11 & 3.1678 & 0.0805\\
img4 & 0.91 & 1.70 & 0.0004 & 0.0009\\
\hline
\end{tabular}
\end{table}

\figura{imagenes/Figuras/FFT/altas_resumen.jpg}{fig:altas}{Pasa altas aplicado a las cuatro imágenes (una por fila): original y resultados con $f_c=0.23$, $0.11$ y $0.91$ (resultado escalado solo para visualizarlo).}
\figura{imagenes/Figuras/FFT/img1_pasa_altas.jpg}{fig:altasdet}{Detalle para img1: por fila $f_c=0.23$, $0.11$ y $0.91$; columnas: original, filtrada (escalada), diferencia y espectro filtrado.}

\paragraph{¿Cuál es la diferencia entre las imágenes originales y las generadas?}
Las imágenes generadas ya no se parecen a las originales: el resultado es un mapa de bordes y detalles finos sobre un fondo gris
uniforme. Las zonas suaves, como el cielo o las superficies lisas, desaparecen y solo quedan los contornos de los objetos. Con
$f_c=0.11$ se conservan más frecuencias intermedias y los bordes tienen más estructura; con $f_c=0.23$ los contornos son más
finos. Con $f_c=0.91$ solo sobreviven las esquinas del espectro y se obtiene un patrón fino casi uniforme, sin contenido
reconocible de la escena y con amplitud casi nula (desviación estándar cercana a $0.001$ en la escala de 0 a 1).

\paragraph{Justificación con el filtro pasa altas.}
La máscara vale 0 dentro del radio $f_c$ y 1 fuera. Como el centro del espectro contiene la frecuencia cero, que representa
el brillo promedio, ese valor se elimina y la imagen resultante tiene media cero, con valores positivos y negativos (por eso
no se recorta a $[0,1]$). Como las dos máscaras suman 1, la imagen original es la suma de su versión pasa bajas y su versión
pasa altas con el mismo $f_c$.

La energía que sobrevive es pequeña: entre 0.6\,\% y 1.5\,\% con $f_c=0.23$, entre 1.7\,\% y 3.8\,\% con $f_c=0.11$ y prácticamente
0\,\% con $f_c=0.91$. En una imagen natural la energía se concentra en las bajas frecuencias, aunque las altas contienen la
información de los bordes. Al bajar $f_c$ el filtro deja pasar más frecuencias (98.1\,\% de los puntos con $f_c=0.11$) y por
tanto más energía y más estructura; al subirlo a $0.91$ queda el 1.7\,\% de los puntos, las frecuencias más altas, y ya no hay
forma de reconocer los objetos. También se observan ondulaciones alrededor de los contornos, por el corte abrupto de la
máscara ideal, igual que en el pasa bajas.



\FloatBarrier
%%%%%%%%%%%%%%%%%%%%%%%%%%%%%%%%%%%%%%%%%%%%%%%%%%%%%%
\subsubsection*{Inciso 3. Blur Detection}

\paragraph{Instrucción.}
Modificar la implementación para que pueda realizar \textit{Blur Detection} y aplicarla a las imágenes proporcionadas.

\paragraph{Respuesta y explicación del código.}
Una imagen nítida tiene bordes, es decir, energía en altas frecuencias; una desenfocada la pierde. Se reutiliza la FFT 2D propia
y se mide la fracción de la energía espectral, sin la componente DC, que cae fuera del radio $r=0.1$:
\begin{equation}
 S=\frac{\sum_{r(u,v)>0.1}|F(u,v)|^2}{\sum_{(u,v)\ne\mathrm{DC}}|F(u,v)|^2}\in[0,1].
\end{equation}
La imagen se declara \textbf{desenfocada} si $S<\tau$ con $\tau=0.02$. Como comprobación independiente se calcula también la
varianza del Laplaciano discreto de cuatro vecinos, método clásico de nitidez (en las tablas se reporta multiplicada por
$255^2$).

\begin{lstlisting}
def puntaje_blur_fft(img, fc_alta=0.1):
    potencia = np.abs(centrar_espectro(fft2_recursiva(img))) ** 2
    r = radio_normalizado(img.shape)
    potencia[img.shape[0] // 2, img.shape[1] // 2] = 0   # se ignora la DC
    return float(potencia[r > fc_alta].sum() / potencia.sum())

def detectar_desenfoque(img, umbral=0.02, fc_alta=0.1):
    puntaje = puntaje_blur_fft(img, fc_alta)
    return {"puntaje_fft": puntaje, "laplaciano": varianza_laplaciano(img),
            "desenfocada": puntaje < umbral}
\end{lstlisting}

\paragraph{Validación con desenfoque artificial.}
Para comprobar que el puntaje responde al desenfoque, cada imagen se desenfocó con un pasa bajas gaussiano
$H(u,v)=e^{-r^2/(2f_c^2)}$ con $f_c=0.30$ (leve), $0.15$ (medio) y $0.05$ (fuerte). Las tablas \ref{tab:blur} y \ref{tab:lap}
y las figuras \ref{fig:blurniveles} y \ref{fig:blurpuntajes} muestran los resultados.

\begin{table}[!htbp]
\centering
\caption{Puntaje $S$ (energía en altas frecuencias) para cada nivel de desenfoque. Umbral $\tau=0.02$.}
\label{tab:blur}
\begin{tabular}{lrrrr}
\hline
\rowcolor{pastelblue}
Imagen & Original & Leve (0.30) & Medio (0.15) & Fuerte (0.05)\\
\hline
img1 & 0.1419 & 0.0911 & 0.0433 & 0.0004\\
img2 & 0.1485 & 0.0898 & 0.0406 & 0.0003\\
img3 & 0.1106 & 0.0732 & 0.0364 & 0.0004\\
img4 & 0.1612 & 0.1002 & 0.0463 & 0.0004\\
\hline
\end{tabular}
\end{table}

\begin{table}[!htbp]
\centering
\caption{Varianza del Laplaciano ($\times255^2$) para los mismos niveles.}
\label{tab:lap}
\begin{tabular}{lrrrr}
\hline
\rowcolor{pastelblue}
Imagen & Original & Leve (0.30) & Medio (0.15) & Fuerte (0.05)\\
\hline
img1 & 1598.8 & 304.7 & 47.3 & 1.4\\
img2 & 1228.8 & 170.6 & 23.5 & 0.7\\
img3 & 880.1 & 151.3 & 26.4 & 0.8\\
img4 & 1266.6 & 200.4 & 26.1 & 0.5\\
\hline
\end{tabular}
\end{table}

\paragraph{Resultados.}
Las cuatro imágenes originales tienen $S$ entre $0.11$ y $0.16$ y se clasifican como \textbf{nítidas}. En todas, el puntaje
decrece estrictamente al aumentar el desenfoque, y la varianza del Laplaciano muestra el mismo orden, lo que respalda la
medida. Con el desenfoque fuerte ($S\approx0.0004$) las cuatro imágenes quedan por debajo del umbral y se detectan como
\textbf{desenfocadas}.

\paragraph{Alcance y limitaciones.}
Con $\tau=0.02$ los desenfoques leve y medio siguen por encima del umbral ($S$ entre $0.036$ y $0.100$) y se clasifican como
nítidos: el detector separa bien la imagen nítida de la muy desenfocada, pero no es sensible a un desenfoque moderado. El
umbral es empírico y el puntaje depende del contenido de la imagen (una escena con pocas texturas tiene menos energía en
altas frecuencias aun estando enfocada), por lo que $\tau$ debe recalibrarse si se cambian las imágenes. No se pretende que
$\tau=0.02$ sea un valor universal.

\figura{imagenes/Figuras/FFT/blur_puntajes.jpg}{fig:blurpuntajes}{
 Puntaje de Blur Detection frente al nivel de desenfoque (escala logarítmica). La línea roja es el umbral $\tau=0.02$.}

\figura{imagenes/Figuras/FFT/blur_niveles_todas.jpg}{fig:blurniveles}{Cada imagen original (fila) y sus versiones con desenfoque leve, medio y fuerte, con su puntaje FFT y la varianza del Laplaciano.}



\FloatBarrier
\subsubsection*{Registro de las consultas a IA (FFT e incisos 1 a 3)}
\evidenciaantonio{FFT y filtros}{fft}{
 La asistencia se utilizó para revisar la implementación frente a las instrucciones y comprobar la convención de frecuencia de corte y las pruebas contra \texttt{np.fft}. El código de la FFT y de los filtros es el que se entrega en el notebook.}
\evidenciaantonio{Inciso 1}{inciso1}{
 Se utilizó para explicar y revisar la interpretación del filtro pasa bajas y su relación con el fenómeno de Gibbs.}
\evidenciaantonio{Inciso 2}{inciso2}{
 Se utilizó para revisar la interpretación del filtro pasa altas y su relación complementaria con el pasa bajas.}
\evidenciaantonio{Inciso 3}{inciso3}{
 Se utilizó para explicar el significado matemático del puntaje y del umbral de decisión. Las cifras de todas las tablas provienen de la ejecución del notebook.}
% AJUSTAR los textos anteriores para que coincidan con las capturas reales de Antonio.

% FIN DEL APARTADO DE ANTONIO PAPADOPOLO CASTELLANOS

\clearpage

%%%%%%%%%%%%%%%%%%%%%%%%%%%%%%%%%%%%%%%%%%%%%%%%%%%%%%
% INICIO DEL APARTADO DE MARIA CLAUDIA LAINFIESTA HERRERA

\subsection*{
    Parte 2 - Redes neuronales convolucionales y comparación
}

\noindent

%%%%%%%%%%%%%%%%%%%%%%%%%%%%%%%%%%%%%%%%%%%%%%%%%%%%%%
\subsubsection*{Inciso 4. Filtro de imagen basado en CNN}

\paragraph{Instrucción.}

Investigar filtros de imagen basados en redes neuronales
convolucionales (CNN) y utilizar un modelo preentrenado para
una tarea de restauración, sin necesidad de entrenar una red
desde cero. En esta implementación se eligió la eliminación
de ruido, también llamada \textit{denoising}.

\paragraph{Preparación común de las imágenes.}

El trabajo se realizó en la sección de CNN de
\path{MAP_final.ipynb}.
Se cargaron \path{img1.jpeg}, \path{img2.png},
\path{img3.jpg} y \path{img4.jpg} desde la carpeta
\path{imagenes/Imagenes de prueba/}.

La función \texttt{cargar\_imagen} del núcleo compartido
convierte cada imagen a escala de grises, la redimensiona a
$256\times256$ píxeles con interpolación Lanczos y normaliza
sus intensidades al intervalo $[0,1]$.

Las dimensiones son potencias de dos, como requiere la FFT
recursiva. Todos los métodos se evalúan contra estas
referencias ya preparadas, no contra los archivos originales
a color o con sus dimensiones de origen.

%%%%%%%%%%%%%%%%%%%%%%%%%%%%%%%%%%%%%%%%%%%%%%%%%%%%%%
\subsubsection*{
    Inciso 4a. Implementación y aplicación del modelo preentrenado
}

\paragraph{Instrucción.}

Agregar la implementación del modelo CNN y utilizarla para
las imágenes adjuntas en el proyecto.

\paragraph{Respuesta y explicación del código.}

Se seleccionó \textbf{DnCNN}, variante
\texttt{dncnn\_gray\_blind}, con pesos publicados en KAIR por
Kai Zhang. El fundamento de DnCNN es estimar el ruido
mediante aprendizaje residual y obtener la restauración
restando esa estimación de la entrada; las referencias del
artículo y del código están al final.

La implementación usada contiene 20 capas convolucionales
con núcleos $3\times3$: una primera capa de un canal a
64 canales, 18 capas intermedias de 64 canales y una última
capa de salida de un canal. Se aplica ReLU después de las
primeras 19 convoluciones.

Se reproduce la configuración de los pesos seleccionados;
en esta clase no se agregan capas de normalización por lotes.

Los pesos se almacenan en
\path{herramientas/model_zoo/dncnn_gray_blind.pth}.
Si el archivo no existe, el código crea la carpeta y lo
descarga de la publicación v1.0 de KAIR.

Se usa \texttt{load\_state\_dict} con
\texttt{strict=True} para comprobar que las capas coincidan
con los pesos. Después, la red se coloca en CPU o GPU,
según disponibilidad, se activa \texttt{eval()} y se
deshabilitan los gradientes. No se realiza entrenamiento
ni actualización de los pesos.

La función \texttt{aplicar\_cnn} verifica que la entrada
sea una matriz válida en $[0,1]$, la convierte en un tensor
con forma $(1,1,256,256)$ y ejecuta el modelo con
\texttt{torch.inference\_mode()}.

La red devuelve

\begin{equation}
    \widehat{x}
    =
    \operatorname{clip}
    \bigl(y-R_{\theta}(y),0,1\bigr),
\end{equation}

donde $y$ es la imagen de entrada y $R_{\theta}(y)$
representa el ruido estimado por las capas convolucionales.
Se verifica la forma y la validez de la salida; el recorte
limita las intensidades al intervalo permitido.

En este inciso se procesan las imágenes cargadas, todavía
sin el ruido artificial de 4b. La comparación visual permite
observar el efecto de DnCNN, pero no prueba por sí sola una
mejora: no se conoce una referencia más limpia que la imagen
original preparada.

La evaluación controlada corresponde a los incisos
siguientes. El notebook también muestra la versión de
PyTorch, el dispositivo, la huella SHA-256 de los pesos y
una tabla de comprobación de dimensiones.

\galeriacnn{4a}{dncnn_comparacion}{
    Aplicación de DnCNN a las cuatro imágenes preparadas
    sin agregar ruido artificial. En cada par, la referencia
    en gris aparece a la izquierda y la salida de DnCNN
    a la derecha.
}

\evidenciacnn{4a}{1}{
    La asistencia se utilizó para organizar la carga del
    modelo preentrenado, su aplicación a las imágenes y las
    comprobaciones de dimensiones. Las capturas deben
    corresponder a la consulta y a la respuesta reales que
    se utilizaron durante este desarrollo.
}

\FloatBarrier

%%%%%%%%%%%%%%%%%%%%%%%%%%%%%%%%%%%%%%%%%%%%%%%%%%%%%%
\subsubsection*{Inciso 4b. Incorporación de ruido artificial}

\paragraph{Instrucción.}

Tomar una imagen original y agregarle ruido artificialmente.
Para ampliar la comparación, se repitió el procedimiento
con las cuatro imágenes del proyecto.

\paragraph{Respuesta y explicación del código.}

Se generó ruido gaussiano con \texttt{default\_rng} de
NumPy, media cero, desviación estándar
$\sigma=25/255\approx0.09804$ y semilla 42.

El procedimiento se expresa como

\begin{equation}
    n_{ij}\sim\mathcal{N}(0,\sigma^2),
    \qquad
    y_{ij}
    =
    \operatorname{clip}(x_{ij}+n_{ij},0,1).
\end{equation}

Aquí $x$ es la referencia en gris y $y$ es la entrada
ruidosa común para FFT y DnCNN.

La semilla se reinicia dentro del ciclo de cada imagen;
por ello, las cuatro imágenes de igual tamaño reciben el
mismo arreglo de ruido antes del recorte. El recorte impide
intensidades negativas o mayores que uno, pero modifica la
distribución del ruido cerca de los extremos del intervalo.

Se crean copias independientes de las referencias y se
comprueba que el resultado conserve las dimensiones y
contenga valores finitos.

Los diccionarios \texttt{originales\_4b} e
\texttt{imagenes\_ruido} conservan las referencias y las
entradas degradadas. También se guardan archivos como
\path{img1_gris.npy} e \path{img1_ruido_gaussiano.npy}
en \path{resultados/ruido_gaussiano/}.

Los incisos 4c, 4d y 5 leen estos mismos archivos: no
generan realizaciones distintas de ruido para cada método.

Visualmente, el ruido introduce granulado en las zonas
uniformes y alrededor de los contornos. El PSNR de estas
entradas se encuentra entre 20.28 y 20.70 dB, y su promedio
es 20.41 dB.

Estos valores sirven como punto de comparación: una
restauración útil debe acercarse más a la referencia que
la entrada ruidosa.

\begin{lstlisting}
#ruido reproducible por imagen y recorte al rango valido
rng = np.random.default_rng(SEMILLA)
ruido = rng.normal(0, SIGMA, original.shape).astype(np.float32)
imagen_ruidosa = np.clip(original + ruido, 0, 1)
\end{lstlisting}

\galeriacnn{4b}{ruido_gaussiano_comparacion}{
    Referencias en escala de grises y sus versiones con
    ruido gaussiano de intensidad $25/255$. Se utiliza
    la misma escala de visualización en ambos paneles.
}

\evidenciacnn{4b}{2}{
    La asistencia se utilizó para preparar ruido
    reproducible, conservar las referencias y guardar las
    entradas comunes del experimento. El código final
    fija la semilla y evita cambiar el ruido entre los
    métodos.
}

\FloatBarrier

%%%%%%%%%%%%%%%%%%%%%%%%%%%%%%%%%%%%%%%%%%%%%%%%%%%%%%
\subsubsection*{Inciso 4c. Limpieza mediante filtros de FFT}

\paragraph{Instrucción.}

Aplicar los filtros de FFT implementados para intentar
limpiar la imagen con ruido.

\paragraph{Respuesta y explicación del código.}

Se importó \texttt{filtro\_pasabajas} de
\path{src/fft_nucleo.py}, sin modificar el núcleo
compartido. Se utilizaron filtros pasa bajas ideales con
frecuencias de corte $f_c=0.15$, $0.22$ y $0.96$.

En este inciso se evalúan esas tres configuraciones;
no se aplican filtros pasa altas.

Para cada imagen se cargan la referencia y la entrada
ruidosa guardadas en 4b. Se comprueba que sean matrices
del mismo tamaño, con valores en $[0,1]$ y dimensiones
que sean potencias de dos.

El filtro calcula la FFT bidimensional propia, centra
el espectro, lo multiplica por una máscara y reconstruye
la imagen con la FFT inversa propia.

La máscara ideal es

\begin{equation}
 H_{f_c}(u,v)
 =
 \begin{cases}
 1, & r(u,v)\leq f_c,\\
 0, & r(u,v)>f_c.
 \end{cases}
\end{equation}

El radio está normalizado desde el centro del espectro:
$r=0$ representa la frecuencia cero y $r=1$ corresponde
a la esquina.

Para una imagen de alto $H$ y ancho $W$, se usa

\begin{equation}
 r(u,v)
 =
 \frac{1}{\sqrt{2}}
 \sqrt{
    \left(\frac{u-H/2}{H/2}\right)^2
    +
    \left(\frac{v-W/2}{W/2}\right)^2
 }.
\end{equation}

La salida se obtiene tomando la parte real de la
transformada inversa y limitándola a $[0,1]$.

No se introduce relleno adicional ni se cambia el tamaño
en esta etapa. Los resultados se conservan en
\texttt{resultados\_fft\_4c} y en
\path{resultados/fft/}, con un archivo \texttt{.npy}
por imagen y por corte.

\begin{lstlisting}
#filtro pasa bajas ideal con la FFT propia
for fc in FC_4C:
    resultado = filtro_pasabajas(ruidosa, fc, tipo="ideal")
\end{lstlisting}

El corte 0.15 reduce más las frecuencias altas, pero
también borra contornos y texturas. El corte 0.22 ofrece
el mejor equilibrio entre reducción de ruido y
conservación de información entre los tres valores probados.

El corte 0.96 deja pasar casi todo el espectro y produce
una salida muy próxima a la entrada ruidosa.

Como el ruido y los detalles finos comparten frecuencias
altas, el pasa bajas no puede eliminarlos de manera
independiente. El corte abrupto de la máscara también
puede producir oscilaciones cerca de los bordes.

\comparacionescnn{4c}{fft_comparacion}{
    Limpieza con FFT de las cuatro imágenes. Cada fila
    muestra la referencia, la entrada ruidosa y las
    salidas de los cortes 0.15, 0.22 y 0.96,
    en ese orden.
}

\evidenciacnn{4c}{3}{
    La asistencia se utilizó para aplicar las funciones
    de FFT existentes a las entradas de 4b, probar los
    tres cortes y guardar las salidas sin alterar el
    núcleo compartido ni regenerar el ruido.
}

\FloatBarrier

%%%%%%%%%%%%%%%%%%%%%%%%%%%%%%%%%%%%%%%%%%%%%%%%%%%%%%
\subsubsection*{Inciso 4d. Limpieza mediante la CNN}

\paragraph{Instrucción.}

Aplicar el filtro basado en IA a la misma imagen con ruido
utilizada para la FFT.

\paragraph{Respuesta y explicación del código.}

Se reutiliza la instancia de DnCNN cargada en 4a y se leen
los mismos archivos \texttt{.npy} de 4b.

Antes de procesar las imágenes, el código comprueba que
el modelo esté disponible y que las referencias y las
entradas ruidosas existan. Se verifican sus dimensiones,
su rango y sus valores finitos.

La entrada se convierte a \texttt{float32}, se organiza
como un tensor de un elemento y un canal y se envía al
dispositivo del modelo.

La inferencia se realiza con la red en modo de evaluación,
sin entrenamiento. En este inciso se llama directamente
al modelo para conservar también la salida anterior al
recorte.

Se guardan dos arreglos por imagen: la salida cruda y la
restaurada limitada a $[0,1]$, en
\path{resultados/dncnn/}. La comparación de 5 utiliza la
salida restaurada.

La tabla \path{resumen_4d.csv} registra las dimensiones
y el archivo de entrada empleado.

\begin{lstlisting}
#inferencia con el modelo cargado en 4a
with torch.inference_mode():
    salida = modelo_dncnn_4a(entrada)

#salida cruda y restauracion limitada al rango valido
salida_cruda = salida[0, 0].cpu().numpy()
resultado = np.clip(salida_cruda, 0, 1)
\end{lstlisting}

La CNN reduce el granulado visible y conserva mejor las
formas y varios contornos que el pasa bajas. Esto no
implica una recuperación perfecta: algunas texturas
pequeñas pueden suavizarse.

A diferencia de una máscara fija en frecuencia, la red
utiliza el contexto espacial de la entrada para estimar
el residuo que debe retirar.

Las imágenes de la figura \ref{fig:cnn-4d} permiten
comparar este efecto con la referencia conocida.

\comparacionesdncnn

\evidenciacnn{4d}{4}{
    La asistencia se utilizó para reutilizar el modelo
    preentrenado, aplicar la inferencia a las mismas
    entradas de FFT y conservar las salidas crudas y
    restauradas para comprobar el procedimiento.
}

\FloatBarrier

%%%%%%%%%%%%%%%%%%%%%%%%%%%%%%%%%%%%%%%%%%%%%%%%%%%%%%
\subsubsection*{Inciso 5. Comparación de los resultados}

\paragraph{Instrucción.}

Indicar cuál filtro funcionó mejor, justificar la respuesta
mediante las imágenes resultantes y explicar cómo la IA
puede superar a filtros clásicos al aprender características
de una imagen limpia.

\paragraph{Procedimiento de comparación.}

Se cargan las referencias de 4b, las tres salidas FFT de 4c
y las salidas DnCNN de 4d. Se incluye la entrada sin filtrar
para medir si la restauración mejora el punto de partida.

La función \texttt{medir\_error} utiliza
\texttt{scikit-image} para calcular MSE, PSNR y SSIM
con rango de datos igual a uno.

Se verifica que cada resultado tenga las mismas dimensiones
que la referencia. No se ajusta el contraste individualmente
ni se retiran bordes o rellenos en este cálculo.

El error cuadrático medio y el PSNR se definen como

\begin{align}
 \operatorname{MSE}(x,\widehat{x})
 &=
 \frac{1}{HW}
 \sum_{i=1}^{H}
 \sum_{j=1}^{W}
 (x_{ij}-\widehat{x}_{ij})^2,
 \\
 \operatorname{PSNR}(x,\widehat{x})
 &=
 10\log_{10}\!
 \left(
    \frac{1}{\operatorname{MSE}(x,\widehat{x})}
 \right).
\end{align}

Un MSE menor indica menor error; un PSNR mayor indica mayor
cercanía bajo ese criterio.

SSIM evalúa la similitud estructural mediante estadísticas
locales de luminancia, contraste y estructura; se interpreta
como mejor una puntuación mayor.

El código emplea su configuración predeterminada y comprueba
que las dimensiones sean de al menos siete píxeles.

Se selecciona la mejor FFT por PSNR y por SSIM de forma
separada usando la referencia conocida. En estos resultados
ambos criterios eligen $f_c=0.22$ para las cuatro imágenes.

Esta selección es retrospectiva entre los tres cortes
probados; no constituye una regla automática para imágenes
nuevas. Las figuras muestran la FFT elegida por PSNR.

\paragraph{Respuesta: ¿cuál filtro funcionó mejor?}

\textbf{DnCNN obtuvo el mejor resultado en las cuatro
imágenes}: presentó mayor PSNR y SSIM que las configuraciones
FFT evaluadas.

La tabla \ref{tab:cnn-psnr} muestra el PSNR de cada imagen
y la tabla \ref{tab:cnn-ssim} muestra la comparación
estructural.

\begin{table}[H]
\centering

\caption{
    PSNR por imagen, en dB.
    La mejor FFT corresponde a $f_c=0.22$.
}
\label{tab:cnn-psnr}

\begin{tabular}{lrrr}
\hline
\rowcolor{pastelblue}
Imagen & Con ruido & Mejor FFT & DnCNN\\
\hline
img1 & 20.70 & 22.69 & \textbf{28.20}\\
img2 & 20.29 & 24.66 & \textbf{29.26}\\
img3 & 20.40 & 25.34 & \textbf{31.53}\\
img4 & 20.28 & 24.31 & \textbf{30.81}\\
\hline
\end{tabular}
\end{table}

\begin{table}[H]
\centering

\caption{
    SSIM por imagen. Un valor mayor representa mayor
    similitud estructural.
}
\label{tab:cnn-ssim}

\begin{tabular}{lrrr}
\hline
\rowcolor{pastelblue}
Imagen & Con ruido & Mejor FFT & DnCNN\\
\hline
img1 & 0.4953 & 0.6451 & \textbf{0.8311}\\
img2 & 0.4394 & 0.7053 & \textbf{0.8673}\\
img3 & 0.4137 & 0.7540 & \textbf{0.9206}\\
img4 & 0.4322 & 0.6481 & \textbf{0.8758}\\
\hline
\end{tabular}
\end{table}

\begin{table}[H]
\centering

\caption{Promedios de las métricas para las cuatro imágenes.}
\label{tab:cnn-promedios}

\begin{tabular}{lrrr}
\hline
\rowcolor{pastelblue}
Método & MSE & PSNR (dB) & SSIM\\
\hline
Sin filtrar
    & 0.009097 & 20.41 & 0.4452\\
FFT pasa bajas (0.15)
    & 0.005433 & 22.80 & 0.6271\\
FFT pasa bajas (0.22)
    & 0.003856 & 24.25 & 0.6881\\
FFT pasa bajas (0.96)
    & 0.009058 & 20.43 & 0.4458\\
DnCNN
    & \textbf{0.001058}
    & \textbf{29.95}
    & \textbf{0.8737}\\
\hline
\end{tabular}
\end{table}

En promedio, DnCNN supera en aproximadamente $5.70$ dB al
mejor corte FFT, calculando la diferencia entre los
promedios de PSNR. También obtiene un MSE menor y un
SSIM mayor.

Los promedios de PSNR se calculan sobre los valores
individuales en dB; no son el PSNR obtenido a partir del
MSE promedio.

Las tablas se obtienen de \path{resumen_5.csv} y
\path{promedios_5.csv}, y el detalle completo está en
\path{metricas_5.csv}, dentro de
\path{resultados/comparacion/}.

\paragraph{¿Por qué puede funcionar mejor una CNN?}

El pasa bajas elimina frecuencias según una máscara fija,
independientemente del contenido local. Así reduce ruido,
pero también descarta componentes que describen bordes y
texturas.

DnCNN utiliza filtros aprendidos previamente y varias capas
que combinan información del entorno de cada píxel para
estimar el ruido.

En la figura \ref{fig:cnn-5} se observa una mayor
conservación de los contornos junto con una reducción del
granulado.

La red no aprende con estas cuatro imágenes: aquí se usan
los pesos que ya fueron entrenados por sus autores.

\paragraph{Alcance de la conclusión.}

La ventaja observada corresponde a estas cuatro referencias
de $256\times256$, al ruido gaussiano con
$\sigma=25/255$ seguido de recorte y a los tres filtros
pasa bajas ideales considerados.

No demuestra que DnCNN supere a cualquier filtro clásico
ni a todos los tipos de ruido.

Además, el mismo patrón aleatorio se reutiliza por imagen
y la mejor FFT se elige conociendo la referencia. Las
métricas cuantifican semejanza con esa referencia, pero no
garantizan que se conserve cada detalle visual.

\comparacionescnn{5}{comparacion}{
    Comparación final. De izquierda a derecha:
    referencia, entrada con ruido, mejor FFT por PSNR
    ($f_c=0.22$ en todos los casos) y DnCNN.
    Las cuatro filas corresponden a img1, img2,
    img3 e img4.
}

\evidenciacnn{5}{5}{
    La asistencia se utilizó para organizar la comparación
    de los métodos, calcular las métricas frente a una
    referencia común y redactar su interpretación.
    Los valores de las tablas corresponden a los
    resultados guardados del experimento, no a cifras
    propuestas por la IA.
}

%%%%%%%%%%%%%%%%%%%%%%%%%%%%%%%%%%%%%%%%%%%%%%%%%%%%%%
\subsubsection*{Conclusión del apartado de CNN}

La comparación permitió relacionar un filtro clásico en
frecuencia con un modelo de restauración aprendido
previamente.

Ambos métodos redujeron el ruido con configuraciones
adecuadas, pero DnCNN obtuvo los mejores valores de las
tres métricas y conservó mejor los contornos en este
experimento.

El almacenamiento de las mismas entradas para ambos
métodos y de las salidas numéricas permitió verificar la
comparación. Las figuras y las métricas respaldan la
conclusión dentro del alcance descrito.

\paragraph{Organización de las salidas.}

Los PNG individuales y las figuras comparativas se guardan
únicamente dentro de \path{imagenes/Figuras/CNN/}, en las
subcarpetas \path{4a/}, \path{4b/}, \path{4c/},
\path{4d/} y \path{5/}.

Los PNG individuales conservan las dimensiones de la imagen
procesada y no incluyen títulos ni ejes. Las figuras
comparativas sí incluyen los títulos para interpretar
el experimento.

Además, se conservan los arreglos numéricos
\texttt{.npy} en \path{resultados/}; estos archivos
permiten calcular las métricas sin depender de la
cuantización o del diseño de los PNG.

% Si hubo mas de una consulta en un inciso,
% agregar todas sus evidencias:
%
% \capturacnn{4a_prompt_02.png}{
%     Segundo prompt del inciso 4a
% }
%
% \capturacnn{4a_respuesta_02.png}{
%     Segunda respuesta del inciso 4a
% }
%
% Agregar debajo la explicacion de como se utilizo
% cada respuesta adicional.

% FIN DEL APARTADO DE MARIA CLAUDIA LAINFIESTA HERRERA

\clearpage


%%%%%%%%%%%%%%%%%%%%%%%%%%%%%%%%%%%%%%%%%%%%%%%%%%%%%%
\logouniversidad

\begin{center}
\Large
\textbf{Referencias}
\end{center}

\begin{itemize}

    \item Google Colab.
    Integración con modelos de IA.
    \url{https://colab.research.google.com/}.

    \item Notas de clase,
    Portal de Matemática VI en GES.

    \item ChatGPT,
    Respuestas generadas por IA,
    ChatGPT, 2026.
    \url{https://chatgpt.com/}.

    \item OpenAI Codex,
    Respuestas generadas por IA (conversación de la Parte 1),
    OpenAI, 2026.

    \item Claude,
    Respuestas generadas por IA,
    Claude, 2026.
    \url{https://claude.com/es}.

    \item Departamento de Matemática Aplicada,
    Matemática 6, FISICC.
    \textit{MAP: Serie de Fourier Truncada y FFT}.
    Guía del proyecto, semestre 6, 2026,
    pp. 4--5.

    \item Zhang, K., Zuo, W., Chen, Y.,
    Meng, D. y Zhang, L.
    \textit{
        Beyond a Gaussian Denoiser:
        Residual Learning of Deep CNN for Image Denoising
    }.
    IEEE Transactions on Image Processing,
    26(7), 3142--3155, 2017.
    \url{https://arxiv.org/abs/1608.03981}.

    \item Zhang, K.
    \textit{
        KAIR: implementación de DnCNN
        y pesos preentrenados
    }.
    \url{https://github.com/cszn/KAIR/blob/master/models/network_dncnn.py}.
    Publicación de pesos v1.0:
    \url{https://github.com/cszn/KAIR/releases/tag/v1.0}.

    \item Cooley, J. W. y Tukey, J. W.
    An algorithm for the machine calculation of complex Fourier series.
    \textit{Mathematics of Computation}, 19(90), 297--301, 1965.

    \item Gonzalez, R. C. y Woods, R. E.
    \textit{Digital Image Processing}. Pearson, 4a. edición, 2018.

    \item Virtanen, P. et al.
    SciPy 1.0: fundamental algorithms for scientific computing in Python.
    \textit{Nature Methods}, 17, 261--272, 2020.

    \item Archivos del experimento:
    \path{MAP_final.ipynb} (notebook entregado para ambas partes),
    \path{resumen_5.csv},
    \path{promedios_5.csv} y
    \path{metricas_5.csv}.

\end{itemize}

\end{document}